% ============================================================
% INSERT AFTER Section 4.6 (Pre-Saturation Spreading Rates)
% and BEFORE Section 5 (Connection to the dm^3 Framework)
% ============================================================

\subsection{Self-Trapping Threshold Gap}
\label{sec:self_trapping}

The spreading-exponent data in Table~\ref{tab:alpha} covers
$\lambda \in [0.5, 2.0]$, the pre-saturation window at $N=2000$.
Extending the fit to the full sweep $\lambda \in \{1, 2, 4, 8, 10\}$
at $N=500$, $T=10^5$ reveals a qualitative transition that the short-time
($T=50$) IPR table does not resolve: the two chains cross distinct
self-trapping thresholds at different values of $\lambda$.

\begin{table}[h]
\centering
\caption{Late-time spreading exponents $\alpha$ from
$\mathrm{IPR}(t) \sim t^{-\alpha}$ fitted over the full $T=10^5$ run at
$N=500$. At $\lambda \geq 4$ the Fibonacci exponent drops to near zero
(self-trapped); the tribonacci exponent remains $\gtrsim 0.1$ through
$\lambda = 10$.}
\label{tab:alpha_full}
\begin{tabular}{r r r l}
\toprule
$\lambda$ & $\alpha_{\rm fib}$ & $\alpha_{\rm trib}$ & Regime \\
\midrule
1  & 0.249 & 0.421 & Both spreading; tribonacci faster \\
2  & 0.119 & 0.243 & Both slowing; tribonacci $\approx 2\times$ \\
4  & 0.040 & 0.184 & Fibonacci near-frozen; tribonacci spreading \\
8  & 0.007 & 0.110 & Fibonacci self-trapped; tribonacci spreading \\
10 & 0.021 & 0.090 & Both suppressed; tribonacci still spreading \\
\bottomrule
\end{tabular}
\end{table}

The data expose a \emph{self-trapping threshold gap}: the Fibonacci chain
crosses into the self-trapped regime ($\alpha \approx 0$) in the interval
$\lambda \in [4, 8]$, while the tribonacci chain has not crossed it by
$\lambda = 10$. This is a qualitative distinction, not merely a
quantitative one.

The short-time result (Section~4.1) reported $\sim\!57\%$ IPR drop for
Fibonacci versus $<\!5\%$ for tribonacci at $\lambda=1.5$, $T=50$: a
large instantaneous differential with both chains still in the spreading
regime. The long-time result is stronger in kind: at $\lambda = 8$ the
Fibonacci chain has essentially ceased spreading
($\alpha_{\rm fib} = 0.007$), while the tribonacci chain is still
propagating the initial perturbation outward
($\alpha_{\rm trib} = 0.110$). The two observations are consistent and
reinforce each other---the $T=50$ snapshot is the early-time face of the
same multifractal resistance that, at long times, shifts the self-trapping
threshold to higher $\lambda$.

\paragraph{Mechanism.}
The standard self-trapping condition for a DNLS site is
$\lambda \langle |\psi_j|^4 \rangle \sim J_{\rm eff}$, where $J_{\rm eff}$
is the effective bandwidth \cite{Eilbeck1985,Kevrekidis2009}.
For a critical multifractal eigenstate the relevant bandwidth is set by
the gap scale $\sim \eta^{-k}$ (tribonacci) or $\sim \varphi^{-k}$
(Fibonacci) at the hierarchical level $k$ corresponding to the
localization length. Because $\eta > \varphi$, the tribonacci gap scale
decays more slowly, providing a higher effective barrier to self-trapping.
This is the same spectral-gap lever identified in the companion criticality
study~\cite{Grossi2026Criticality}, where the critical fission strength
$\lambda_c(n)$ is governed by the spectral gap $\Delta_n = \rho_n -
|\rho_n^{(2)}|$ of the substitution transfer matrix. The present result is
the nonlinear-dynamics counterpart: a higher spectral gap lifts the
self-trapping threshold.

\paragraph{Limitations.}
The $\alpha$ values in Table~\ref{tab:alpha_full} are fitted over a finite
time window at finite $N=500$; saturation effects are not fully excluded,
and the $\lambda=10$ Fibonacci uptick ($\alpha_{\rm fib} = 0.021 >
\alpha_{\rm fib}|_{\lambda=8} = 0.007$) may reflect a secondary oscillatory
mode rather than genuine residual spreading. We therefore interpret the
self-trapping threshold as lying in $\lambda \in [4, 8]$ for Fibonacci and
$\lambda > 10$ for tribonacci, with the exact values to be resolved by
a finer $\lambda$-scan (listed as an open question in Section~6).

% ============================================================
% REPLACE item 4 in the Open Questions enumerate (Section 6)
% with the following:
% ============================================================
%
% \item \textbf{Self-trapping threshold} (\emph{partially closed}).
% Section~\ref{sec:self_trapping} establishes that the Fibonacci
% self-trapping threshold lies in $\lambda \in [4, 8]$ and the
% tribonacci threshold is $> 10$ at $N=500$, $T=10^5$.
% A fine $\lambda$-scan in the interval $[2, 12]$ at fixed $N=500$,
% $T=10^5$ would bracket both thresholds precisely. Preliminary
% tetrabonacci ($n=4$) results suggest a further upward shift consistent
% with the spectral-gap-controls-threshold picture.
%
% ============================================================
% ADD this bibitem to \begin{thebibliography} if not already present:
% ============================================================
%
% \bibitem{Grossi2026Criticality}
% P.~Nogueira~Grossi,
% ``Criticality Thresholds in One-Dimensional Multiplying Media
% with $n$-Bonacci Aperiodic Modulation:
% Spectral Gap Control of $k_{\rm eff}$ in
% Substitution-Sequence Diffusion Operators.''
% G6~LLC. Zenodo:
% \href{https://doi.org/10.5281/zenodo.20077205}{10.5281/zenodo.20077205}
% (2026).
%
% ============================================================
% REPLACE the final paragraph of Section 7 (Conclusion) with:
% ============================================================
%
% The self-trapping analysis (Section~\ref{sec:self_trapping}) adds a
% third, qualitative layer to the differential-robustness picture.
% At $T=50$ the Fibonacci state has lost $\sim\!57\%$ of its linear IPR
% while the tribonacci state retains $>95\%$ (differential in magnitude).
% At $T=10^5$ the Fibonacci spreading exponent collapses to
% $\alpha \approx 0$ for $\lambda \geq 8$ while the tribonacci exponent
% remains $\gtrsim 0.1$ through $\lambda=10$ (differential in regime).
% These are two faces of the same mechanism: the tribonacci spectral gap
% $\Delta_3 = \eta - |\eta^{(2)}| > \Delta_2 = \varphi - |\varphi^{(2)}|$
% raises the effective self-trapping barrier.  The companion criticality
% paper~\cite{Grossi2026Criticality} shows the identical lever operating
% in a classical neutron-diffusion setting, where $\lambda_c(n) \approx
% 0.958\,\Delta_n + 0.107$ across $n = 2, \ldots, 5$ with $r = 0.989$.
% The convergence of two independent physical models on the same
% spectral-gap control parameter constitutes the primary finding of this
% pair of papers.
